\chapter{Universo matematico e controparte non standard}\stepcounter{capitoloeccolo}

Dopo aver mostrato la costruzione di $ \starR $ come ultraprodotto, si utilizza una costruzione simile, generalizzandola ad un \emph{universo matematico}. Questo consentirà di superare parzialmente la limitazione della scelta del linguaggio $ \mathcal{L} $, e quindi di parlare degli oggetti matematici ``comuni'' e della loro controparte non standard, pur senza uscire dal formalismo logico. 

%TODO: Mostrare il perché voglio fare quest costruzione:
% Noto che, poiché l'obiettivo di studio di questa tesi sono i manifolds, allora non è sufficiente studiare semplicemente R*, bensì c'è bisogno di una struttura più grande. Si introdurrà quindi il concetto di sovrastruttura di un insieme X, e se ne costruirà *V(X). 
% Si riconoscerà poi un'immersione elementare tra R* (come definita nel chap precedente) e R* all'interno di *V(R). Vedi retro tesi
% All'interno di V(RR), poi, si identificherà il manifold M.
